% ID: SEL-D0077-OB-E03
% ESTADO: cotejada. Segunda pasada visual y borrador compilado revisados.
% Fuente: Selectividad/sel_2008_matematicas_aplicadas/Mat CCSS 5 - Examen y criterios de correción.pdf
% SHA-256: 6b6f35fe4872415542ee008d454a2502926c6ed0fffe87639ed13bed50759c6b
% Página: 2; opción B; ejercicio original 3.

\noindent\underline{Parte I}

Se dispone de los siguientes datos sobre el equipamiento de los hogares de una ciudad:

En el $60\%$ de los hogares se puede ver la TDT (Televisión Digital Terrestre) y el $70\%$ de los hogares dispone de ordenador. De entre los hogares que disponen de ordenador, el $80\%$ puede ver la TDT.
\begin{enumerate}[label=\alph*)]
\item \textbf{(1 punto)} \textquestiondown{}Son sucesos independientes ``disponer de ordenador'' y ``poder ver la TDT''?
\item \textbf{(1 punto)} \textquestiondown{}Qué porcentaje de hogares no disponen de ordenador ni pueden ver la TDT?
\end{enumerate}
\par\noindent\begin{minipage}{\linewidth}
\noindent\underline{Parte II}

\textbf{(2 puntos)} En un centro de anillamiento de aves se ha detectado que en una muestra de 250 ejemplares de una especie, 60 son portadoras de una bacteria.

Obtenga un intervalo de confianza, al $97\%$, para la proporción de aves de esa especie que son portadoras de la bacteria.
\end{minipage}\par
